\chapter{模型与算法}

\section{模型看什么、输出什么}

正式部署模型属于视觉—语言—动作模型：它把三路相机画面、当前双臂状态和任务指令放在一起理解，然后直接生成双臂关节与夹爪的未来动作。该类模型可以把预训练得到的视觉与语言知识迁移到机器人操作；本项目采用的 $\pi_{0.5}$ 系列强调开放场景下的长程操作与多来源共同训练\cite{pi05_2025}。不过，本报告的结构和数字全部以项目实际保存模型和训练记录为准，不把论文中的通用能力自动当成本项目已实现能力。

\figfull[.98\textwidth]{model_dataflow.pdf}{正式部署模型的数据流。三路视野、14 个机器人状态量和任务指令共同形成动作；每次输出未来 50 个控制时刻。}

\begin{table}[H]
\centering
\caption{正式模型输入与输出}
\begin{tabularx}{\textwidth}{>{\bfseries}p{31mm}p{34mm}X}
\toprule
项目 & 形状/规模 & 作用\\
\midrule
顶部画面 & $224\times224\times3$ & 识别桌面布局、瓶子候选和两臂整体关系\\
左腕画面 & $224\times224\times3$ & 观察左夹爪、瓶身和遮挡区域\\
右腕画面 & $224\times224\times3$ & 观察瓶口、瓶盖、右夹爪和旋转接触\\
机器人状态 & 14 个离散化状态量 & 告诉模型双臂与夹爪当前在哪里\\
任务指令 & 最长 200 个语言单位 & 说明要完成的任务步骤\\
动作输出 & $50\times14$ & 未来 50 个控制时刻的双臂关节和夹爪动作\\
\bottomrule
\end{tabularx}
\end{table}

\section{整体结构}

实际保存模型包含图像编码部分、语言与任务理解部分以及动作生成部分。图像先被划分为小块并编码；三路图像、语言和离散化机器人状态进入统一时序模型；动作专家根据这些信息产生连续动作片段。模型内部使用 18 层与动作相关的注意力记录，真实导出结果覆盖 50 个动作查询和每路 $16\times16$ 的图像网格。\ev{E09,E12}

\begin{table}[H]
\centering
\caption{正式部署模型配置}
\begin{tabularx}{\textwidth}{>{\bfseries}p{38mm}p{35mm}X}
\toprule
项目 & 核验值 & 说明\\
\midrule
模型类型 & $\pi_{0.5}$ 视觉—语言—动作模型 & 使用连续动作生成头\\
保存参数量 & 838,358,468 & 51 个参数叶节点；仅保存模型参数\\
数值精度 & bfloat16 & 训练计算精度\\
图像输入 & 三路 $224\times224$ & 不使用底部相机\\
状态/真实动作 & 14 / 14 维 & 内部动作容器为 32 维，输出时只取机器人实际维度\\
动作片段长度 & 50 & 对应约 1 秒控制计划\\
生成修正次数 & 10 & 从随机动作开始逐步变为可执行动作\\
训练方式 & 全参数微调 & 当次配置未冻结参数\\
预训练来源 & $\pi_{0.5}$ 基础模型 & 具体预训练数据不属于本项目自采数据\\
指数滑动平均 & 训练配置为 0.99 & 部署保存文件中未发现独立平均参数副本\\
\bottomrule
\end{tabularx}
\end{table}

\section{动作学习目标}

训练时，设真实示教动作片段为 $\mathbf{a}$，随机噪声为
$\boldsymbol{\epsilon}\sim\mathcal{N}(0,I)$，随机时间
$t=0.001+0.999\,\operatorname{Beta}(1.5,1)$。模型看到的中间动作是
\[
\mathbf{x}_t=t\boldsymbol{\epsilon}+(1-t)\mathbf{a},
\]
正确的“修正方向”为
\[
\mathbf{u}_t=\boldsymbol{\epsilon}-\mathbf{a}.
\]
模型预测 $\mathbf{v}_\theta(\mathbf{x}_t,t,\mathbf{o})$，其中
$\mathbf{o}$ 代表三路画面、状态和任务指令。实际训练目标为
\[
\mathcal{L}_{\mathrm{flow}}
=\mathbb{E}\!\left[
\frac{1}{D}\left\|
\mathbf{v}_\theta(\mathbf{x}_t,t,\mathbf{o})-\mathbf{u}_t
\right\|_2^2\right],
\]
$D=32$ 是内部动作容器宽度；机器人真实使用其中 14 维。\ev{E06}

\plain{训练时先把正确动作“弄乱”，再让模型学会应该往哪个方向修正，才能把随机动作一步步变回像人类示教的动作。}

推理时从随机动作 $\mathbf{x}_1$ 开始，用 10 个等长小步积分到 $\mathbf{x}_0$：
\[
\mathbf{x}_{t-\Delta t}
=\mathbf{x}_{t}-\Delta t\,\mathbf{v}_\theta(\mathbf{x}_t,t,\mathbf{o}),
\qquad \Delta t=0.1.
\]
最终得到 50 步动作片段。该公式与实际项目实现的符号方向一致。

\section{训练配置与资源}

\begin{table}[H]
\centering
\caption{正式部署模型训练配置}
\begin{tabularx}{\textwidth}{>{\bfseries}p{37mm}p{38mm}X}
\toprule
项目 & 核验值 & 解释\\
\midrule
批量规模 & 256 & 每次更新使用的示教样本量\\
并行计算 & 4 张 H200 & 分布式全参数训练\\
数据读取进程 & 64 & 提高多仓库与视频读取吞吐\\
随机种子 & 42 & 只有一个种子，无法计算种子波动\\
学习率 & 峰值 $2.5\times10^{-5}$ & 先 10,000 步升温，再衰减至 $2.5\times10^{-6}$\\
优化器 & AdamW 系列 & 梯度裁剪上限 1.0\\
保存间隔 & 每 1,000 步 & 现场采用第 19,000 步保存版本\\
历史覆盖 & 0--59,990 步 & 记录跨度约 51.4 小时\\
训练/验证/测试 & 只有训练 & 当次记录没有验证和测试指标\\
\bottomrule
\end{tabularx}
\end{table}

训练平台中的原始配置写的是 40,000 步，但实际历史继续到 59,990 步。这一冲突在报告中保留：可以确认训练继续发生，不能仅凭现有记录断言 40,000 步之后是否改变过数据权重。现场部署版本是 19,000 步，不是损失最低点或训练终点。\ev{E08}

\section{注意力记录能解释到什么程度}

\figfull[\textwidth]{attention_real_example.png}{真实运行中的三路注意力样例。彩色热区表示生成动作时更集中查看的位置；该记录没有成功/失败标签。}

\figfull[.82\textwidth]{attention_camera_share.pdf}{8,223 份真实注意力记录的三路视野份额。各份额是归一化比较；腕部近景整体高于顶部总览。}

9 份运行清单共包含 8,223 个真实样本，格式一致、无解析失败。各运行的平均份额范围为：顶部总览 22.9\%--25.9\%，左腕 32.9\%--36.4\%，右腕 37.7\%--43.4\%。首次记录后的导出开销中位数约 40--55 毫秒，95 分位约 41--57 毫秒。\ev{E12}

\limitbox{注意力只能说明模型在生成动作时更集中查看哪些画面区域。没有瓶盖存在性、瓶身方向和成功/失败标签，也没有遮挡热区后的动作变化对照，因此不能用热图直接证明空拧或倒抓的因果原因。}

\section{部署保存模型审计}

正式部署保存模型是只含参数的分块格式，目录标记为第 19,000 步；包含 51 个参数叶节点、838,358,468 个参数、14 维状态和动作的均值、标准差与第 1/99 百分位统计。保存文件没有优化器状态，也没有独立的滑动平均参数。\ev{E09}

保存模型的形状与当前部署结构一致，能够读取元数据。需要保守说明的是：保存文件内部没有另一个独立字段再次写明“19,000”，因此步数是由目录、运行记录和部署配置联合确认，而不是单凭文件名猜测。

\begin{table}[H]
\centering
\caption{正式保存模型核查结果}
\begin{tabularx}{\textwidth}{>{\bfseries}p{41mm}p{31mm}X}
\toprule
核查项 & 结果 & 解释\\
\midrule
是否能读取元数据 & 能 & 参数结构和统计信息可只读解析\\
保存格式 & 仅模型参数 & 不包含可继续训练所需的完整优化器状态\\
参数节点/总参数 & 51 / 838,358,468 & 与正式模型结构一致\\
训练步数 & 联合确认19,000 & 由部署配置、目录和训练谱系交叉确认\\
归一化统计 & 存在 & 状态与动作各14维，含均值、标准差和分位数\\
独立平均参数 & 未发现 & 训练设置有平均系数，但保存文件没有单独副本\\
当前代码兼容性 & 元数据与形状一致 & 没有发现明确形状冲突\\
\bottomrule
\end{tabularx}
\end{table}
